% GENERATED FILE. Edit canonical lessons, then run npm run generate:books.
% canonical-source-hash: fnv1a64:b45903f59793762b
% canonical-lessons: TE-C198-tyusanu, TE-C198-intarvyu, TE-C198-siphtu, TE-C198-yunipharam, TE-C198-sibbandi, TE-R198-first-pass-words-from-chapters-193-195, TE-R198-second-pass-words-from-chapters-196-198

\chapter{Tuition, Interview, Work shift, Uniform, Staff}
\label{ch:te-tuition-interview-work-shift-uniform-staff}
\hlchaptermodality{198}

\begin{chapteropening}
\textbf{By the end of this chapter:} \emph{I can say tuition, an interview, a work shift, a uniform and staff.}

Complete the last lesson of chapter 198: I can say tuition, an interview, a work shift, a uniform and staff.
\end{chapteropening}

\section[\texorpdfstring{\te{ట్యూషను} (ṭyūṣanu)}{ట్యూషను (ṭyūṣanu)}]{\te{ట్యూషను} (ṭyūṣanu) — tuition, coaching}

\label{lesson:TE-C198-tyusanu}



\begin{quote}
Before the new one: say the Telugu for a solution, then the Telugu for an opinion.
\end{quote}

\subsection*{You'll want to know: \te{ట్యూషను}}
\textbf{\te{ట్యూషను}} — \emph{ṭyūṣanu} — \textquotedblleft{}tuition, coaching\textquotedblright{}.

\begin{grammarlens}[title={What you've built}]
One more word for this chapter.
\end{grammarlens}

\subsection*{Your turn}
\begin{itemize}
\raggedright
  \item \emph{Say it:} \emph{ṭyūṣanu}
  \item \emph{Say it:} \emph{ṭyūṣanu}, once more
  \item \emph{Say it:} say \emph{ṭyūṣanu}
  \item \emph{Recall:} say the Telugu for a solution, then the Telugu for an opinion, then say \emph{ṭyūṣanu} again
\end{itemize}

\begin{tcolorbox}[breakable,colback=teal!4,colframe=teal!35!black,title={Before you move on}]
What does \te{ట్యూషను} mean? (\textquotedblleft{}Tuition, coaching\textquotedblright{}.) Say it once more.
\end{tcolorbox}
\section[\texorpdfstring{\te{ఇంటర్వ్యూ} (iṇṭarvyū)}{ఇంటర్వ్యూ (iṇṭarvyū)}]{\te{ఇంటర్వ్యూ} (iṇṭarvyū) — an interview}

\label{lesson:TE-C198-intarvyu}



\begin{quote}
Before the new one: say the Telugu for an opinion, then the Telugu for tuition.
\end{quote}

\subsection*{You'll want to know: \te{ఇంటర్వ్యూ}}
\textbf{\te{ఇంటర్వ్యూ}} — \emph{iṇṭarvyū} — \textquotedblleft{}an interview\textquotedblright{}.

\begin{grammarlens}[title={What you've built}]
One more word for this chapter.
\end{grammarlens}

\subsection*{Your turn}
\begin{itemize}
\raggedright
  \item \emph{Say it:} \emph{iṇṭarvyū}
  \item \emph{Say it:} \emph{iṇṭarvyū}, once more
  \item \emph{Say it:} say \emph{iṇṭarvyū}
  \item \emph{Recall:} say the Telugu for an opinion, then the Telugu for tuition, then say \emph{iṇṭarvyū} again
\end{itemize}

\begin{tcolorbox}[breakable,colback=teal!4,colframe=teal!35!black,title={Before you move on}]
What does \te{ఇంటర్వ్యూ} mean? (\textquotedblleft{}An interview\textquotedblright{}.) Say it once more.
\end{tcolorbox}
\section[\texorpdfstring{\te{షిఫ్టు} (ṣiphṭu)}{షిఫ్టు (ṣiphṭu)}]{\te{షిఫ్టు} (ṣiphṭu) — a work shift}

\label{lesson:TE-C198-siphtu}



\begin{quote}
Before the new one: say the Telugu for tuition, then the Telugu for an interview.
\end{quote}

\subsection*{You'll want to know: \te{షిఫ్టు}}
\textbf{\te{షిఫ్టు}} — \emph{ṣiphṭu} — \textquotedblleft{}a work shift\textquotedblright{}.

\begin{grammarlens}[title={What you've built}]
One more word for this chapter.
\end{grammarlens}

\subsection*{Your turn}
\begin{itemize}
\raggedright
  \item \emph{Say it:} \emph{ṣiphṭu}
  \item \emph{Say it:} \emph{ṣiphṭu}, once more
  \item \emph{Say it:} say \emph{ṣiphṭu}
  \item \emph{Recall:} say the Telugu for tuition, then the Telugu for an interview, then say \emph{ṣiphṭu} again
\end{itemize}

\begin{tcolorbox}[breakable,colback=teal!4,colframe=teal!35!black,title={Before you move on}]
What does \te{షిఫ్టు} mean? (\textquotedblleft{}A work shift\textquotedblright{}.) Say it once more.
\end{tcolorbox}
\section[\texorpdfstring{\te{యూనిఫారం} (yūniphāraṁ)}{యూనిఫారం (yūniphāraṁ)}]{\te{యూనిఫారం} (yūniphāraṁ) — a uniform}

\label{lesson:TE-C198-yunipharam}



\begin{quote}
Before the new one: say the Telugu for an interview, then the Telugu for a work shift.
\end{quote}

\subsection*{You'll want to know: \te{యూనిఫారం}}
\textbf{\te{యూనిఫారం}} — \emph{yūniphāraṁ} — \textquotedblleft{}a uniform\textquotedblright{}.

\begin{grammarlens}[title={What you've built}]
One more word for this chapter.
\end{grammarlens}

\subsection*{Your turn}
\begin{itemize}
\raggedright
  \item \emph{Say it:} \emph{yūniphāraṁ}
  \item \emph{Say it:} \emph{yūniphāraṁ}, once more
  \item \emph{Say it:} say \emph{yūniphāraṁ}
  \item \emph{Recall:} say the Telugu for an interview, then the Telugu for a work shift, then say \emph{yūniphāraṁ} again
\end{itemize}

\begin{tcolorbox}[breakable,colback=teal!4,colframe=teal!35!black,title={Before you move on}]
What does \te{యూనిఫారం} mean? (\textquotedblleft{}A uniform\textquotedblright{}.) Say it once more.
\end{tcolorbox}
\section[\texorpdfstring{\te{సిబ్బంది} (sibbandi)}{సిబ్బంది (sibbandi)}]{\te{సిబ్బంది} (sibbandi) — staff}

\label{lesson:TE-C198-sibbandi}



\begin{quote}
Before the new one: say the Telugu for a work shift, then the Telugu for a uniform.
\end{quote}

\subsection*{You'll want to know: \te{సిబ్బంది}}
\textbf{\te{సిబ్బంది}} — \emph{sibbandi} — \textquotedblleft{}staff\textquotedblright{}.

\begin{grammarlens}[title={What you've built}]
One more word for this chapter.
\end{grammarlens}

\subsection*{Your turn}
\begin{itemize}
\raggedright
  \item \emph{Say it:} \emph{sibbandi}
  \item \emph{Say it:} \emph{sibbandi}, once more
  \item \emph{Say it:} say \emph{sibbandi}
  \item \emph{Recall:} say the Telugu for a work shift, then the Telugu for a uniform, then say \emph{sibbandi} again
\end{itemize}

\begin{tcolorbox}[breakable,colback=teal!4,colframe=teal!35!black,title={Before you move on}]
What does \te{సిబ్బంది} mean? (\textquotedblleft{}Staff\textquotedblright{}.) Say it once more.
\end{tcolorbox}
\section[(dialogue)]{Words from chapters 193-195 — a first pass}

\label{lesson:TE-R198-first-pass-words-from-chapters-193-195}



\begin{quote}
Say the words for a uniform and staff.
\end{quote}

\subsection*{Your turn}
Say these aloud:

\begin{itemize}
\raggedright
  \item a line, a circle, a square, a triangle, a wheel
  \item housework, a domestic helper, a tailor, a trader, a coin
  \item an officer, a clerk, a manager, a colleague, a meeting
\end{itemize}

\begin{tcolorbox}[breakable,colback=teal!4,colframe=teal!35!black,title={Before you move on}]
A line? (\textbf{\te{గీత}}, \emph{gīta}.) A circle? (\textbf{\te{వృత్తం}}, \emph{vṛttaṁ}.) A square? (\textbf{\te{చతురస్రం}}, \emph{caturasraṁ}.) A triangle? (\textbf{\te{త్రిభుజం}}, \emph{tribhujaṁ}.) A wheel? (\textbf{\te{చక్రం}}, \emph{cakraṁ}.) Housework? (\textbf{\te{ఇంటిపని}}, \emph{iṇṭipani}.) A domestic helper? (\textbf{\te{పనిమనిషి}}, \emph{panimaniṣi}.) A tailor? (\textbf{\te{దర్జీ}}, \emph{darjī}.) A trader? (\textbf{\te{వ్యాపారి}}, \emph{vyāpāri}.) A coin? (\textbf{\te{నాణెం}}, \emph{nāṇeṁ}.) An officer? (\textbf{\te{అధికారి}}, \emph{adhikāri}.) A clerk? (\textbf{\te{గుమాస్తా}}, \emph{gumāstā}.) A manager? (\textbf{\te{మేనేజరు}}, \emph{mēnējaru}.) A colleague? (\textbf{\te{సహోద్యోగి}}, \emph{sahōdyōgi}.) A meeting? (\textbf{\te{సమావేశం}}, \emph{samāvēśaṁ}.)
\end{tcolorbox}
\section[(dialogue)]{Words from chapters 196-198 — a second pass}

\label{lesson:TE-R198-second-pass-words-from-chapters-196-198}



\begin{quote}
Say the words for a fee and a fine.
\end{quote}

\subsection*{Your turn}
Say these aloud:

\begin{itemize}
\raggedright
  \item a fee, a fine, a deadline, a responsibility, an application
  \item a carpenter, a writer, a problem, a solution, an opinion
  \item tuition, an interview, a work shift, a uniform, staff
\end{itemize}

\begin{tcolorbox}[breakable,colback=teal!4,colframe=teal!35!black,title={Before you move on}]
A fee? (\textbf{\te{ఫీజు}}, \emph{phīju}.) A fine? (\textbf{\te{జరిమానా}}, \emph{jarimānā}.) A deadline? (\textbf{\te{గడువు}}, \emph{gaḍuvu}.) A responsibility? (\textbf{\te{బాధ్యత}}, \emph{bādhyata}.) An application? (\textbf{\te{దరఖాస్తు}}, \emph{darakhāstu}.) A carpenter? (\textbf{\te{వడ్రంగి}}, \emph{vaḍraṅgi}.) A writer? (\textbf{\te{రచయిత}}, \emph{racayita}.) A problem? (\textbf{\te{సమస్య}}, \emph{samasya}.) A solution? (\textbf{\te{పరిష్కారం}}, \emph{pariṣkāraṁ}.) An opinion? (\textbf{\te{అభిప్రాయం}}, \emph{abhiprāyaṁ}.) Tuition? (\textbf{\te{ట్యూషను}}, \emph{ṭyūṣanu}.) An interview? (\textbf{\te{ఇంటర్వ్యూ}}, \emph{iṇṭarvyū}.) A work shift? (\textbf{\te{షిఫ్టు}}, \emph{ṣiphṭu}.) A uniform? (\textbf{\te{యూనిఫారం}}, \emph{yūniphāraṁ}.) Staff? (\textbf{\te{సిబ్బంది}}, \emph{sibbandi}.)
\end{tcolorbox}
